%% kapa: exam=2 round=36 year=2025 period=3 subject=compensation_law subject_ko=감정평가 및 보상법규 seq=4 number=4 points=10
%% source: Q-Net 감정평가사 기출문제 (cst00309) artlSeq=5239325 file=raw/1008/5239325/2025년 제36회 감정평가사 2차 감정평가 및 보상법규.pdf
%% conversion: kordoc 4.24.51 -> tools/gen2.py (auto); review=proofread-gpt-sol; flags=none
\begin{longquestion}{【문제 4】}
A 등은 감정평가법인을 설립하여 국토교통부장관의 설립인가를 받았다. 설립인가 자체의 하자는 없지만, 감정평가법인의 사원에 관한 ｢감정평가 및 감정평가사에 관한 법률｣의 요건을 충족하지 못하여 설립행위에 문제가 있음을 확인한 경쟁업체의 감정평가사 甲은 소송을 통하여 이를 다투고자 한다. 이러한 감정평가법인의 설립행위와 인가행위의 관계를 설명하고, 甲은 어떠한 소송에서 어떠한 하자를 다툴 수 있는지 검토하시오. (10점) \par
\end{longquestion}
